\chapter{Continuité, valeurs intermédiaires et équations}
\section*{Exercice 13.1 -- Existence et unicité}
La fonction est continue sur $\R$, sa dérivée $3x^2+2$ est strictement positive. Elle est donc strictement croissante. Ses valeurs en $0$ et $1$ sont $-2$ et $1$ : le TVI assure une racine entre les deux, la croissance son unicité globale.

\section*{Exercice 13.2 -- Dichotomie}
L'erreur est au plus $1/2^{n+1}$. Avec $n=9$, elle vaut au plus $1/1024<10^{-3}$ ; avec $n=8$, la borne $1/512$ ne suffit pas. Ne pas confondre longueur de l'intervalle et erreur du milieu.

\section*{Exercice 13.3 -- Prolongement}
Le quotient est le taux d'accroissement de l'exponentielle en $0$, donc sa limite vaut $\exp'(0)=1$. Définir la valeur en zéro égale à $1$ réalise le prolongement.

\section*{Exercice 13.4 -- TVI sans monotonie}
Elle possède deux solutions $\pm1/2$. La continuité ne garantit pas l'unicité. En outre les valeurs aux deux extrémités sont égales à $1$ ; pour utiliser directement le TVI, séparer les intervalles $[-1,0]$ et $[0,1]$.

\section*{Exercice 13.5 -- Image d’une suite}
La fonction $x\mapsto\sqrt{1+x^2}$ est continue sur $\R$ : polynôme positif puis racine continue. L'image de la suite converge donc vers $\sqrt5$. La justification est la continuité, pas une substitution purement formelle.

\section*{Exercice 13.6 -- Additivité continue}
L'additivité impose $f(q)=3q$ pour tout rationnel $q$. Choisir des rationnels $q_n\to\sqrt2$ ; la continuité donne $f(\sqrt2)=\lim f(q_n)=3\sqrt2$.
